\subsection{Zwölf Karten Poker}
Die folgende Pokervariante ist von mir selbst entworfen worden, als eine Zwischenstufe der Komplexität zwischen Leduc Hold'em und Rhode Island Hold'em.
Die 12 Karten Variante verwendet ein Deck mit zwölf Karten: J, Q, K und A in drei Suits (s, h, d).

Es gibt drei Runden.
In der ersten Runde erhält jeder Spieler eine private Karte.
In der zweiten und dritten Runde wird jeweils eine öffentliche Board Karte aufgedeckt.

Es gibt ein Maximum von zwei Bets pro Runde.
Die Setzgrößen betragen 2 Chips in der ersten Runde, 4 Chips in der zweiten Runde und 8 Chips in der dritten Runde.
Zu Beginn zahlen beide Spieler einen Ante von 1 Chip in den Pot.

Das Spiel endet nach der dritten Runde mit einem Showdown.
Die Hand besteht aus der Kombination der privaten Karte des Spielers und den beiden öffentlichen Board Karten.
Die Hand Rankings sind: Three of a Kind schlägt ein Paar, ein Paar schlägt High Card.


